% front-matter/preface.tex
%
% LNCS front-matter: Preface
% \input-ed by front-matter/front-matter.tex; do NOT add \begin{document} here.
%
% ============================================================
% EDITORIAL NOTES -- remove this block before final submission
%
% 1. KEYNOTE ABSTRACTS
%    Abstracts of keynote talks that do not have an accompanying full paper belong
%    in the Front Matter, not in the main proceedings body. Place them after the
%    Organization pages and before the first research paper, following Section 4.2
%    of the Springer editor guidelines. Do NOT include keynote-only entries in
%    tracking/paper-tracking.md or submission/toc.csv.
%
% 2. POST-EVENT REVISION SUBTITLE
%    If authors were given the opportunity to revise their papers after the event,
%    the volume subtitle must reflect this (e.g., "Revised Selected Papers"). If
%    papers are published as submitted, no qualifier is added. Confirm the correct
%    subtitle with Springer when requesting the volume number.
%    TODO: record here whether post-event revision was offered -- Yes / No
%
% 3. VOLUME EDITORS = PROGRAM COMMITTEE CHAIRS (DEFAULT)
%    Springer's default is that the volume editors are the Program Committee Chairs.
%    If a different arrangement is intended, contact Ronan Nugent
%    (ronan.nugent@springer.com) before the volume number is assigned.
%
% 4. FRONT-MATTER TEMPLATES
%    Official LaTeX and Word templates for the LNCS front matter (preface,
%    organization pages, title page) are available from the volume editors'
%    website at http://bit.ly/2HaWmko. Use those templates if the structure
%    required by Springer differs from this stub.
%
% 5. WRITING STYLE
%    Write the preface in American English, past tense. The preface is typically
%    one to two pages in the final LNCS volume. Avoid first-person singular; use
%    "we" for the editorial team.
% ============================================================

\title{Preface}
\author{}
\institute{}
\maketitle

%------------------------------------------------------------------
% Event Information (pre-populated -- verify before submission)
%------------------------------------------------------------------
VERIF-AI 2026 was held in Villebrumier, France, on 8--11~March 2026,
organized by the LASER Foundation. The workshop brought together
researchers and practitioners working on the verification, validation,
and assurance of AI-based systems, with the proceedings published in
the Springer Lecture Notes in Computer Science series.

%------------------------------------------------------------------
\section*{Scope and Aims}
%------------------------------------------------------------------
% TODO: Describe the scientific scope of the workshop: the core research
% questions it addressed, the community it serves, and what distinguishes
% VERIF-AI from related venues. Approximately 2--3 paragraphs.
% Write in American English, past tense.

%------------------------------------------------------------------
\section*{Review Process}
%------------------------------------------------------------------
% TODO: Describe the submission and review process: who was invited to
% submit, the number of submissions received, the review model
% (single-blind / double-blind), the number of reviewers per paper,
% and any discussion or meta-review phase. Write in American English,
% past tense.

%------------------------------------------------------------------
\section*{Acceptance Statistics}
%------------------------------------------------------------------
% TODO: Report the acceptance statistics. Example sentence to adapt:
% "The program committee received N full paper submissions. After a
% rigorous review process, 25 papers were selected for inclusion in
% these proceedings, representing an acceptance rate of approximately
% X\%."

%------------------------------------------------------------------
\section*{Acknowledgements}
%------------------------------------------------------------------
% TODO: Thank the Program Committee members, additional reviewers, the
% local organizing team, the LASER Foundation, and any sponsors or
% funding agencies that supported the workshop. Write in American
% English, past tense.

%------------------------------------------------------------------
% Volume Editors signature block
%------------------------------------------------------------------
% TODO: Replace the placeholders below with the full names, affiliations,
% and countries of the volume editors (one block per editor).
% Default: volume editors = Program Committee Chairs (see note 3 above).

\bigskip
\noindent
March 2026 \hfill [PLACEHOLDER -- Volume Editor Name]\\
\phantom{March 2026 \hfill}[PLACEHOLDER -- Affiliation, Country]

\smallskip
\noindent
\phantom{March 2026 \hfill}[PLACEHOLDER -- Volume Editor Name]\\
\phantom{March 2026 \hfill}[PLACEHOLDER -- Affiliation, Country]
